\documentclass[../../main.tex]{subfiles}
\graphicspath{{\subfix{../../image/}}}

\begin{document}

\section*{第一章小结}\label{section:ch01-小结}

本章中介绍了一类新的现象——随机现象，这是一种普遍存在的现象。在大量随机现象中存在着统计规律性，概率论便是研究随机现象的数量规律的一门数学学科。

``事件''与``概率''是概率论中最基本的两个概念，我们在公理化结构中严格地定义了这两个概念。

为了使读者清楚地理解事件与概率的直观意义，我们采用由具体到抽象，由简单到复杂，由特殊到一般的方式分别介绍了频率、古典概型、几何概率，并从中归纳出事件与概率的本质特征，为公理化定义作准备，这种讲法基本上与概率概念的历史发展平行。

事件的运算及概率的性质是本章的基本内容，也是学习以后各章的必要基础，务必牢固掌握。

我们较完整地研究了古典概型，并介绍了它在产品抽样检查中的应用。对于古典概型的讨论有助于对概率论基本概念的直观理解，而且在以后讨论更一般的情况时，也常以它为特例加以考察。古典概型中概率的计算有较高的技巧性，读者应该掌握一些最基本的计算方法。

几何概率是很有启发性的一类问题，不过它的严格表述只有用到第三章的一些概念才能做到。

搞清频率与概率的关系是十分重要的一个课题，今后我们将一再回到这个问题上来。

书中我们着重把现代概率论公理化结构作为一个历史发展过程来描述，有兴趣的读者不妨试着进行相反相成的另一项工作：用公理化结构来概括古典概型、几何概率等特殊模型。

尽管公理化结构是对抽象样本空间给定的，但应用中我们还是最关心 Euclid 空间，在这种场合，Borel 点集是一个重要的概念，在第三章中，它将起关键作用。

\end{document}